../cictikz/src/cictikz/data/tex/cictikz_preamble.tex

% Why a source follower exists: a photodiode dumps its charge on a tiny
% 1 fF node. Buffered by the follower, the sense node keeps its full
% signal swing while the follower drives the big 1 pF load.

\begin{circuitikz}[american, thick, transform shape, circuit ee IEC]

  ../cictikz/src/cictikz/data/tex/cictikz_lib.tex

  \def\ysig{2.4}    % the sense node rail

  % photodiode into the sense node
  \draw (0,0) \vground;
  %- compact filled diode, house weight
  \draw (0,0) -- (0,{\ysig/2-0.25});
  \fill (-0.2,{\ysig/2-0.25}) -- (0.2,{\ysig/2-0.25}) -- (0,{\ysig/2+0.25}) -- cycle;
  \draw (-0.2,{\ysig/2+0.25}) -- (0.2,{\ysig/2+0.25});
  \draw (0,{\ysig/2+0.25}) -- (0,\ysig);
  %- the two arrows that make it a photodiode rather than a diode
  \draw[->] (1.15,0.55) -- (0.45,0.95);
  \draw[->] (1.15,0.95) -- (0.45,1.35);
  %- Polarity matters and is easy to get backwards. The reset switch
  %- charges the sense node to 3.0 V, so the cathode has to sit on the
  %- node: the junction is reverse biased, and light discharges it. Turn
  %- the diode around and it is forward biased at 3 V, clamping the node
  %- near 0.6 V, and there is no pixel left. The photocurrent therefore
  %- runs the other way to the symbol's triangle, down and out of the
  %- node, which is where the minus sign in dV comes from.
  \draw[red, ->] (-0.65,1.55) -- (-0.65,0.85);
  \node[red, anchor=east, scale=0.85] at (-0.7,1.2) {$i_{ph}$};
  %- the transfer gate: the n region is fully depleted, so the electron
  %- packet it collected moves across in one go
  \draw (0,\ysig) to [nos] (1.6,\ysig);
  \node[anchor=south, scale=0.85] at (0.8,{\ysig+0.45}) {TX};

  % the 1 fF sense capacitor
  \draw (2.4,0) -- (2.4,0.4) \vcapacitor{1fF} -- (2.4,2.4);
  \draw (2.4,0) \vground;
  \fill (2.4,\ysig) circle (0.075);
  \draw (1.6,\ysig) -- (3.6,\ysig);

  % reset switch to the supply
  \draw (3.6,\ysig) to [nos] (3.6,{\ysig+1.5});
  \node[anchor=west, scale=0.85] at (3.75,{\ysig+0.75}) {RST};
  \draw (3.6,{\ysig+1.5}) \vsupply;
  \node[anchor=south] at (3.6,{\ysig+1.85}) {3.0V};
  \fill (3.6,\ysig) circle (0.075);
  \node[anchor=south, color=blue] at (2.9,{\ysig+0.6}) {$\Delta V$};
  \draw (3.6,\ysig) -- (4.4,\ysig);

  % the source follower
  \draw (4.4,\ysig) node[nmos, anchor=G] (M1) {};
  \draw (M1.D) -- ++(0,0.4) coordinate (vddtap) \vsupply;
  \node[anchor=south] at ([yshift=14pt]vddtap) {3.0V};

  % bias current and the big load
  \draw (M1.S) -- ++(0,-0.4) coordinate (osrc);
  \draw (osrc) to [I, l_=5$\mu$A] ++(0,-1.8) \vground;
  \draw (osrc) to [short,*-] ++(1.6,0) coordinate (oc);
  \draw (oc) ++(0,-1.8) coordinate (ocb) \vground;
  \draw (ocb) \vcapacitor{1pF} -- (oc);

\end{circuitikz}
\end{document}
